% Part III -- Reusable and Reliable Programs
% Chapter 6: Functions and Reusable Code

\h{Functions and Reusable Code}

By this point, you can store data, make decisions, and repeat instructions.
Functions let you give a useful group of instructions a name. You can then run
that group whenever it is needed, with different inputs each time.

\begin{NoteBox}
\textbf{Prerequisites.} You should understand variables, lists, dictionaries,
conditionals, and loops.
\end{NoteBox}

\begin{tcolorbox}[title=Learning outcomes,colframe=orange,colback=white,arc=2mm]
After completing this chapter, you should be able to:
\begin{itemize}
    \item define and call a function;
    \item distinguish parameters from arguments;
    \item return a value and use it in another calculation;
    \item use positional, keyword, and default arguments;
    \item explain local scope; and
    \item combine small functions into a readable program.
\end{itemize}
\end{tcolorbox}

\hh{Why functions help}

Repeated code is harder to correct. If the same calculation appears in five
places, a change must be made five times. A function keeps the calculation in
one named place.

\code{c06_repeated_calculation}{python}{
    lobby_area_m2 = 6.0 * 8.0
    office_area_m2 = 4.0 * 5.5

    print(lobby_area_m2)
    print(office_area_m2)
}{The same calculation written twice}

The next example moves the shared rule into a function.

\hh{Define and call a function}

Use \c{def} to define a function. The indented body contains the instructions
that belong to it.

\code{c06_first_function}{python}{
    def show_welcome():
        print("Welcome to Chapter 6")

    show_welcome()
    show_welcome()
}{Define a function and call it twice}

Defining a function does not run its body. The body runs when Python reaches a
call such as \c{show_welcome()}.

\begin{NoteBox}
\textbf{Stop and predict.} How many lines appear if the two calls are removed?
The definition remains, but nothing calls it.
\end{NoteBox}

\hh{Give a function input}

A \emph{parameter} is a name in the function definition. An \emph{argument} is
the value supplied in a call.

\code{c06_one_parameter}{python}{
    def show_room(room_name):
        print(f"Room: {room_name}")

    show_room("Lobby")
    show_room("Office")
}{Pass one argument to a parameter}

In this example, \c{room_name} is the parameter. \c{"Lobby"} and
\c{"Office"} are arguments.

\hh{Return a result}

Use \c{return} when the caller needs a value. A returned value can be stored,
printed, compared, or passed to another function.

\code{c06_return_area}{python}{
    def rectangle_area_m2(width_m, length_m):
        area_m2 = width_m * length_m
        return area_m2

    lobby_area_m2 = rectangle_area_m2(6.0, 8.0)
    office_area_m2 = rectangle_area_m2(4.0, 5.5)

    print(lobby_area_m2)
    print(office_area_m2)
}{Return a calculation result}

When Python reaches \c{return}, the function ends and sends the value back to
the call.

\hh{Return and print have different purposes}

Printing displays information. Returning gives a value back to the caller.

\code{c06_print_and_return}{python}{
    def display_area(width_m, length_m):
        print(width_m * length_m)

    def calculate_area(width_m, length_m):
        return width_m * length_m

    displayed = display_area(3.0, 4.0)
    calculated = calculate_area(3.0, 4.0)

    print(f"Displayed returned: {displayed}")
    print(f"Calculated returned: {calculated}")
}{Compare printing with returning}

The first function displays \c{12.0} but returns \c{None}. The second returns
\c{12.0}, so the program can reuse it.

\hh{Use several parameters}

Parameters are matched with arguments from left to right unless keywords are
used.

\code{c06_multiple_parameters}{python}{
    def volume_m3(width_m, length_m, height_m):
        return width_m * length_m * height_m

    result_m3 = volume_m3(5.0, 8.0, 3.2)
    print(f"Volume: {result_m3:.2f} m3")
}{Use three positional arguments}

Choose parameter names that make their roles clear. A reader should not need to
guess what each number represents.

\hh{Use keyword arguments}

Keywords connect each value to a parameter name. They are helpful when a
function has several similar inputs.

\code{c06_keyword_arguments}{python}{
    def volume_m3(width_m, length_m, height_m):
        return width_m * length_m * height_m

    result_m3 = volume_m3(
        length_m=8.0,
        height_m=3.2,
        width_m=5.0
    )

    print(f"Volume: {result_m3:.2f} m3")
}{Call a function with keyword arguments}

Keyword arguments may appear in a different order because their names show how
they match.

\hh{Provide a default value}

A default makes an argument optional. Parameters without defaults must come
before parameters with defaults.

\code{c06_default_argument}{python}{
    def area_with_waste(area_m2, waste_rate=0.05):
        return area_m2 * (1 + waste_rate)

    print(area_with_waste(100.0))
    print(area_with_waste(100.0, waste_rate=0.10))
}{Use and override a default argument}

Defaults should represent a sensible common case. A caller can still provide a
different value.

\hh{Keep names local}

A name created inside a function is local to that call. Code outside the
function cannot use it directly.

\code{c06_local_scope}{python}{
    def calculate_perimeter_m(width_m, length_m):
        perimeter_m = 2 * (width_m + length_m)
        return perimeter_m

    result_m = calculate_perimeter_m(4.0, 7.0)
    print(result_m)
}{Use a local variable and return its value}

The name \c{perimeter_m} exists inside the function. The caller stores the
returned value under the separate name \c{result_m}.

Avoid changing global variables inside a function. Inputs and return values
make the function's behavior easier to see and test.

\hh{Use a guard clause}

A guard clause handles a special case near the top of a function. It prevents
the main calculation from becoming deeply nested.

\code{c06_guard_clause}{python}{
    def average(values):
        if len(values) == 0:
            return None

        return sum(values) / len(values)

    print(average([10.0, 20.0, 30.0]))
    print(average([]))
}{Return early for an empty collection}

Chapter 8 shows how a function can raise an exception when invalid input should
be reported as an error.

\hh{Combine small functions}

One function can call another. This is called composition.

\code{c06_composition}{python}{
    def rectangle_area_m2(width_m, length_m):
        return width_m * length_m

    def material_cost(area_m2, price_per_m2):
        return area_m2 * price_per_m2

    area_m2 = rectangle_area_m2(6.0, 8.0)
    cost = material_cost(area_m2, 12.50)

    print(f"Area: {area_m2:.2f} m2")
    print(f"Cost: ${cost:.2f}")
}{Compose two focused functions}

Each function has one clear responsibility. The short main section shows how
the pieces fit together.

\hh{Document a function}

A docstring is a string placed immediately below the function header. It
briefly explains the function's purpose.

\code{c06_docstring}{python}{
    def celsius_to_fahrenheit(celsius):
        """Return a Celsius temperature converted to Fahrenheit."""
        return celsius * 9 / 5 + 32

    print(celsius_to_fahrenheit(20))
}{Add a short docstring}

Use \c{help(celsius_to_fahrenheit)} in the REPL to display the docstring.

\hh{Common mistakes}

\begin{itemize}
    \item \textbf{Defining but not calling}: the body does not run until a call occurs.
    \item \textbf{Forgetting the colon}: a function header ends with a colon.
    \item \textbf{Incorrect indentation}: every body instruction must be indented.
    \item \textbf{Printing instead of returning}: the caller receives \c{None}.
    \item \textbf{Using an unclear name}: prefer a verb such as
          \c{calculate_area} over a vague name such as \c{do_it}.
    \item \textbf{Writing one very long function}: divide it into focused steps.
\end{itemize}

\hh{Practice ladder}

\hhh{Level 0 -- Read and predict}

\code{c06_practice_level0}{python}{
    def double(value):
        return value * 2

    first = double(4)
    second = double(first)
    print(second)
}{Trace two function calls}

\hhh{Level 1 -- Modify}

Change \c{double} to \c{triple}. Update the name, calculation, and calls.

\hhh{Level 2 -- Complete}

Complete a function that returns the perimeter of a rectangle.

\code{c06_practice_level2}{python}{
    def rectangle_perimeter_m(width_m, length_m):
        return 2 * (width_m + length_m)

    print(rectangle_perimeter_m(5.0, 8.0))
}{A completed perimeter function}

\hhh{Level 3 -- Apply}

Write \c{circle_area_m2(radius_m)}. Use \c{3.14159} for pi. Return the area,
then display the result to two decimal places for a radius of 4 meters.

\textbf{Hints}
\begin{itemize}
    \item \textbf{Hint 1}: Area equals pi multiplied by radius squared.
    \item \textbf{Hint 2}: Use \c{radius_m ** 2}.
    \item \textbf{Hint 3}: Format the returned result outside the function.
\end{itemize}

\hhh{Level 4 -- Challenge}

Write \c{room_report(name, width_m, length_m)}. It should use a separate area
function and return one formatted sentence containing the name and area.

\hh{Practice solutions}

\textbf{Level 0:} The output is \c{16}.

\textbf{Level 2:} The output is \c{26.0}.

\textbf{Level 3}

\code{c06_solution_level3}{python}{
    def circle_area_m2(radius_m):
        return 3.14159 * radius_m ** 2

    area_m2 = circle_area_m2(4.0)
    print(f"Area: {area_m2:.2f} m2")
}{One solution to Level 3}

\textbf{Level 4}

\code{c06_solution_level4}{python}{
    def rectangle_area_m2(width_m, length_m):
        return width_m * length_m

    def room_report(name, width_m, length_m):
        area_m2 = rectangle_area_m2(width_m, length_m)
        return f"{name}: {area_m2:.2f} m2"

    print(room_report("Lobby", 6.0, 8.0))
}{One solution to Level 4}

\hh{Chapter checkpoint}

\begin{enumerate}
    \item Which keyword begins a function definition?
    \item What is the difference between a parameter and an argument?
    \item Does defining a function run its body?
    \item What does \c{return} do?
    \item What does a function return when it has no \c{return} statement?
    \item Why can keyword arguments improve readability?
    \item Where does a local variable exist?
    \item What is a guard clause?
    \item Why are several small functions often clearer than one long function?
\end{enumerate}

\hh{Checkpoint answers}

\begin{enumerate}
    \item \c{def}.
    \item A parameter is a name in the definition; an argument is a supplied value.
    \item No. A call runs it.
    \item It ends the function and sends a value to the caller.
    \item \c{None}.
    \item They show which value belongs to which parameter.
    \item Inside the function call where it was created.
    \item An early check and return for a special case.
    \item Each small function can have one clear responsibility and be tested separately.
\end{enumerate}

\begin{tcolorbox}[title=Chapter summary,colframe=orange,colback=white,arc=2mm]
You can define functions, pass arguments, return results, use defaults and
keywords, and combine focused functions. The next chapter introduces useful
functions and modules that Python already provides.
\end{tcolorbox}

